\documentclass[10pt,a4paper]{amsart}
\usepackage{amsmath,amssymb,amsthm,graphicx}
\usepackage{verbatim}
\newenvironment{reference}{\comment}{\endcomment}
%\newenvironment{reference}{}{}
\newenvironment{slogan}{\comment}{\endcomment}
\newenvironment{history}{\comment}{\endcomment}

% For commutative diagrams we use Xy-pic
\usepackage[all]{xy}

% We use 2cell for 2-commutative diagrams.
\xyoption{2cell}
% Theorem environments.
%
\theoremstyle{plain}
\newtheorem{theorem}[subsection]{Theorem}
\newtheorem{proposition}[subsection]{Proposition}
\newtheorem{lemma}[subsection]{Lemma}

\theoremstyle{definition}
\newtheorem{definition}[subsection]{Definition}
\newtheorem{example}[subsection]{Example}
\newtheorem{exercise}[subsection]{Exercise}
\newtheorem{situation}[subsection]{Situation}

\theoremstyle{remark}
\newtheorem{remark}[subsection]{Remark}
\newtheorem{remarks}[subsection]{Remarks}

\numberwithin{equation}{subsection}

% Macros
%
\def\lim{\mathop{\mathrm{lim}}\nolimits}
\def\colim{\mathop{\mathrm{colim}}\nolimits}
\def\Spec{\mathop{\mathrm{Spec}}}
\def\Hom{\mathop{\mathrm{Hom}}\nolimits}
\def\Ext{\mathop{\mathrm{Ext}}\nolimits}
\def\SheafHom{\mathop{\mathcal{H}\!\mathit{om}}\nolimits}
\def\SheafExt{\mathop{\mathcal{E}\!\mathit{xt}}\nolimits}
\def\Sch{\mathit{Sch}}
\def\Mor{\mathop{\mathrm{Mor}}\nolimits}
\def\Ob{\mathop{\mathrm{Ob}}\nolimits}
\def\Sh{\mathop{\mathit{Sh}}\nolimits}
\def\NL{\mathop{N\!L}\nolimits}
\def\CH{\mathop{\mathrm{CH}}\nolimits}
\def\proetale{{pro\text{-}\acute{e}tale}}
\def\etale{{\acute{e}tale}}
\def\QCoh{\mathit{QCoh}}
\def\Ker{\mathop{\mathrm{Ker}}}
\def\Im{\mathop{\mathrm{Im}}}
\def\Coker{\mathop{\mathrm{Coker}}}
\def\Coim{\mathop{\mathrm{Coim}}}

% Boxtimes
%
\DeclareMathSymbol{\boxtimes}{\mathbin}{AMSa}{"02}

%
% Macros for moduli stacks/spaces
%
\def\QCohstack{\mathcal{QC}\!\mathit{oh}}
\def\Cohstack{\mathcal{C}\!\mathit{oh}}
\def\Spacesstack{\mathcal{S}\!\mathit{paces}}
\def\Quotfunctor{\mathrm{Quot}}
\def\Hilbfunctor{\mathrm{Hilb}}
\def\Curvesstack{\mathcal{C}\!\mathit{urves}}
\def\Polarizedstack{\mathcal{P}\!\mathit{olarized}}
\def\Complexesstack{\mathcal{C}\!\mathit{omplexes}}
% \Pic is the operator that assigns to X its picard group, usage \Pic(X)
% \Picardstack_{X/B} denotes the Picard stack of X over B
% \Picardfunctor_{X/B} denotes the Picard functor of X over B
\def\Pic{\mathop{\mathrm{Pic}}\nolimits}
\def\Picardstack{\mathcal{P}\!\mathit{ic}}
\def\Picardfunctor{\mathrm{Pic}}
\def\Deformationcategory{\mathcal{D}\!\mathit{ef}}
\usepackage{cleveref}
\crefname{originalreference}{original reference}{original references}
\Crefname{originalreference}{Original reference}{Original references}
\makeatletter
\newlabel{lemma-local-isomorphism-fully-faithful}{{096L}{source}{}{sourcepointer-1}{}}
\newlabel{lemma-local-isomorphism-fully-faithful@cref}{{[originalreference][0][]096L}{[1][0][]source}{}{}{}}
\makeatother
\crefname{originalreference}{original reference}{original references}
\Crefname{originalreference}{Original reference}{Original references}
\makeatletter
\newlabel{lemma-localization}{{096V}{source}{}{sourcepointer-2}{}}
\newlabel{lemma-localization@cref}{{[originalreference][0][]096V}{[1][0][]source}{}{}{}}
\makeatother
\crefname{originalreference}{original reference}{original references}
\Crefname{originalreference}{Original reference}{Original references}
\makeatletter
\newlabel{lemma-make-w-local}{{0975}{source}{}{sourcepointer-3}{}}
\newlabel{lemma-make-w-local@cref}{{[originalreference][0][]0975}{[1][0][]source}{}{}{}}
\makeatother
\crefname{originallemma}{original lemma}{original lemmas}
\Crefname{originallemma}{Original lemma}{Original lemmas}
\makeatletter
\newlabel{topology-lemma-bijective-map}{{08YE}{source}{}{sourcepointer-4}{}}
\newlabel{topology-lemma-bijective-map@cref}{{[originallemma][0][]08YE}{[1][0][]source}{}{}{}}
\makeatother
\title{Localizing a ring along a closed profinite subset produces a w-local ring with that set as its closed points}
\author{The Stacks Project Authors}
\date{stacks-project:proetale.tex:local-snapshot-20261012; GNU FDL 1.2 or later}
\begin{document}
\maketitle
\noindent\textit{Editorial note.} All delivered mathematical text is the author's own, copied from the pinned Stacks Project chapter file proetale.tex; no statement or proof step was rewritten, reordered or removed.\par\medskip
\noindent\textit{Editorial note.} The counted proof is the author's complete argument: the identification of the inverse image T with V(I) through the homeomorphism of profinite spaces, the construction of the localization along T with its w-locality and closed points, and the isomorphism of the quotients through faithfulness of local isomorphisms.\par\medskip
\noindent\textit{Editorial note.} The label lemma-localize-along-closed-profinite names the counted statement, and the proof environment immediately following it in the source is delivered verbatim as its proof.\par\medskip
\noindent\textit{Editorial note.} The counted statement is carried by the author's own lemma environment in the chapter source.\par\medskip
\noindent\textit{Editorial note.} Pre-open gate: python operations/agent-authoring/tools/precard\_gate.py --source proetale.tex --label lemma-localize-along-closed-profinite --start 1011 --end 1031 --local-id 28005 returned RESULT=OK (L28005 free, occupancy FREE, public interval clear). Candidate taken from sources/stacks-proetale/candidates-halfB.json.\par\medskip
\noindent\textit{Editorial note.} The statement and proof are copied from the pinned chapter file; the cited inputs are kept as external pointers with their canonical Stacks tags.\par\medskip
\noindent\textit{Allowed standard prerequisites (editorial).}\begin{itemize}
\item The w-localization and its set of closed points. Same chapter, the construction of the w-localization A -> A\_w with its closed point set Z mapping bijectively onto the spectrum of A. Supplies the w-localization A\_w and the closed subset Z used to cut out the profinite subset.
\item Homeomorphism of a bijective map of profinite spaces. Topology, the lemma that a continuous bijection from a profinite space to a compact Hausdorff space is a homeomorphism. Supplies the homeomorphism T -> V(I) for the inverse image T of V(I) in Z.
\item The localization construction and faithfulness of local isomorphisms. Same chapter, the localization construction attached to a locally closed subscheme and the lemma that a local isomorphism which is a homeomorphism on spaces is an isomorphism. Supply the ring B = (A\_w)\_T\textasciicircum{}tilde and the conclusion that A/I -> B/IB is an isomorphism.
\end{itemize}\medskip
\ifdefined\hypertarget\hypertarget{sourcepointer-1}{}\fi
\noindent\textit{Source reference.} Original reference 096L: sources/stacks-proetale/pinned-original-source.tex, source lines 603--603. Prerequisite pointer invoked at proof line(s) 1030; it names a standard fact of the same chapter that the counted argument uses, and every step of the counted proof built on it is retained verbatim in the delivered body.\par\medskip
\ifdefined\hypertarget\hypertarget{sourcepointer-2}{}\fi
\noindent\textit{Source reference.} Original reference 096V: sources/stacks-proetale/pinned-original-source.tex, source lines 710--710. Prerequisite pointer invoked at proof line(s) 1025; it names a standard fact of the same chapter that the counted argument uses, and every step of the counted proof built on it is retained verbatim in the delivered body.\par\medskip
\ifdefined\hypertarget\hypertarget{sourcepointer-3}{}\fi
\noindent\textit{Source reference.} Original reference 0975: sources/stacks-proetale/pinned-original-source.tex, source lines 844--844. Prerequisite pointer invoked at proof line(s) 1021; it names a standard fact of the same chapter that the counted argument uses, and every step of the counted proof built on it is retained verbatim in the delivered body.\par\medskip
\ifdefined\hypertarget\hypertarget{sourcepointer-4}{}\fi
\noindent\textit{Source reference.} Original lemma 08YE: sources/stacks-chow/tags.txt, source lines 11337--11337. Prerequisite pointer invoked at proof line(s) 1024; it names a standard fact of a sibling chapter that the counted argument uses, and every step of the counted proof built on it is retained verbatim in the delivered body.\par\medskip
\begin{lemma}
\label{lemma-localize-along-closed-profinite}
Let $A$ be a ring. Let $V(I) \subset \Spec(A)$ be a closed subset
which is a profinite topological space. Then there exists an
ind-Zariski ring map $A \to B$ such that $\Spec(B)$ is w-local,
the set of closed points is $V(IB)$, and $A/I \cong B/IB$.
\end{lemma}

\begin{proof}
Let $A \to A_w$ and $Z \subset Y = \Spec(A_w)$ as in
Lemma \ref{lemma-make-w-local}.
Let $T \subset Z$ be the inverse image of $V(I)$.
Then $T \to V(I)$ is a homeomorphism by
Topology, Lemma \ref{topology-lemma-bijective-map}.
Let $B = (A_w)_T^\sim$, see Lemma \ref{lemma-localization}.
It is clear that $B$ is w-local with closed points $V(IB)$.
The ring map $A/I \to B/IB$ is ind-Zariski
and induces a homeomorphism on underlying
topological spaces. Hence it is an isomorphism by
Lemma \ref{lemma-local-isomorphism-fully-faithful}.
\end{proof}

\end{document}
